\documentclass[11pt]{article}
\usepackage[margin=1in]{geometry}
\usepackage[T1]{fontenc}
\usepackage{longtable}
\usepackage{array}
\begin{document}
\# FLCR-37 - Plasma, Energy, And Traversal Calibration

**Series:** Formalizing LCR, Unifying Standard Models
**Artifact:** formal paper source
**Status:** first-pass rich rewrite; requires final citation and build pass.
**Claim posture:** maximal local claim posture with explicit lane boundaries.

\subsection*{Abstract}

This paper formalizes plasma and energy translation from traversal ledgers and carrier horizons. The operative object is energy-plasma translation. The core result is that internal traversal cost and carrier thresholds can be translated into plasma/energy hypotheses only with unit conversion and measured observables. The paper also defines how this result routes forward: FLCR-37 composes FLCR-24, FLCR-25, and FLCR-27 with calibration boundaries. Its main residue is explicit: laboratory plasma behavior, joule conversion, and universal energy bounds require external data.

\subsection*{Keywords}

energy-plasma translation; LCR; receipt; claim lane; normal form.

\subsection*{External Reader Orientation}

An outside reviewer should read this paper as a translation and comparator paper. Its local object is **Plasma, Energy, And Traversal Calibration**. The paper's immediate contribution is: Maps plasma/energy claims to unit, witness, and calibration requirements.

The nearest external anchors are standard mathematical physics definitions, cited Standard Model targets, calibrated constants or data where available, and the LCR-native papers that provide the source-side object. LCR terms are therefore internal labels for the construction unless a row explicitly imports an external theorem, citation, dataset, or calibration target.

It may close the external target definition and the internal translation grammar, but it does not claim measured physical identity unless a calibration/comparator row is present.

\subsection*{Reviewer Compression}

**Core object.** Plasma, Energy, And Traversal Calibration.

**Primary result.** This paper contributes a controlled mapping between a cited external target and a cited LCR-native source object.

**Primary non-result.** This paper does not claim measured physical identity unless the comparator/calibration row is attached.

**Review strategy.** Evaluate the formal claim rows first, then inspect the receipt/citation/calibration rows that support each claim. Appendices provide routing and implementation context; they should not be read as stronger than their lane permits.

\subsection*{Peer Reviewer Assessment}

Reviewer assumption: the reviewer has been briefed on the FLCR structure, claim lanes, CAM/crystal routing, forge/action surfaces, and the distinction between internal formal results and external physical calibration.

**Recommendation.** major revision, technically promising.

**Review score vector.** orientation=2, claim\_granularity=2, evidence\_visibility=2, falsifier\_visibility=2, journal\_apparatus=1, external\_target\_boundary=2.

**Formal claim rows reviewed.** 3.

\subsubsection*{Reviewer Strength Finding}

The manuscript correctly treats Standard Model language as an external target surface and keeps translation claims separate from measured physical identity.

\subsubsection*{Major Revision Requests}

\noindent{}-- Convert the strongest proof sketch into numbered definitions, lemmas, and propositions with explicit dependencies.\\
\noindent{}-- Replace placeholder source classes with final citation keys, receipt hashes, or dataset identifiers wherever possible.\\
\noindent{}-- Move repeated pass-derived material into a compact evidence appendix and keep the main body focused on the paper-local argument.\\
\noindent{}-- Add a comparator table mapping each external target object to the exact LCR source object, invariant, calibration status, and falsifier.\\

\subsubsection*{Minor Revision Requests}

\noindent{}-- Normalize theorem capitalization and punctuation.\\
\noindent{}-- Add final figure/table captions where the text currently describes diagrams schematically.\\
\noindent{}-- Ensure every acronym and local term appears in the paper-local glossary.\\
\noindent{}-- State scale, scheme, and units for any physics-facing numerical comparator.\\

\subsubsection*{Acceptance Blockers}

\noindent{}-- Measured physical identity cannot be accepted without comparator/calibration rows.\\
\noindent{}-- The mapping from external target to LCR source must be explicit enough that another reviewer could reject or reproduce it.\\


\subsection*{1. Contribution And Scope}

\noindent{}-- Defines energy-plasma translation as a first-class FLCR object.\\
\noindent{}-- States the local result: internal traversal cost and carrier thresholds can be translated into plasma/energy hypotheses only with unit conversion and measured observables.\\
\noindent{}-- Routes downstream use through claim lanes rather than inherited prose: FLCR-37 composes FLCR-24, FLCR-25, and FLCR-27 with calibration boundaries.\\
\noindent{}-- Preserves the residue boundary: laboratory plasma behavior, joule conversion, and universal energy bounds require external data.\\

This paper belongs to the Standard Model translation band. Standard Model language names an external target surface, not an automatic LCR identity. A translation claim is admissible only when the cited standard object, the LCR-native source object, the mapping rule, the evidence lane, and the falsifier or calibration boundary are all visible.

\subsection*{2. Source Routing}

Primary routed sources:

\noindent{}-- `25\_Energetic\_Traversal\_Maps.md`\\
\noindent{}-- `26\_Z\_Pinch\_and\_Shear\_Horizon.md`\\
\noindent{}-- `29\_Monster\_Universal\_Energy\_Bound\_Hypotheses.md`\\

Cross-corpus evidence classes:

\noindent{}-- `CAM\_CLAIM\_100\_SCORE\_AUDIT.md`\\
\noindent{}-- `NSIT\_TOOL\_CLOSURE\_MATRIX.md`\\
\noindent{}-- `NSIT\_REASONED\_CLOSURE\_PASS.md`\\
\noindent{}-- `PAPER\_REWORKS\_CRYSTAL\_PROJECTION.json`\\
\noindent{}-- standards, glue, window, and node databases where applicable\\
\noindent{}-- notebook-derived review prompts and orbital source manifests\\

\subsection*{3. Definitions}

\noindent{}-- **Observable.** A measured quantity with units and acquisition context.\\
\noindent{}-- **Traversal calibration.** The binding between internal ledger cost and physical energy units.\\
\noindent{}-- **Receipt boundary.** The exact scope in which the paper's result can be replayed or consumed.\\
\noindent{}-- **Promotion route.** The evidence path required before a stronger downstream claim can use this result.\\

\subsection*{4. Formal Claims}

\noindent\texttt{| Claim | Lane | Statement |}\\
\noindent\texttt{|-------|------|-----------|}\\
\noindent\texttt{| Theorem 37.1 | `external\_calibration\_result` | internal traversal cost and carrier thresholds can be translated into plasma/energy hypotheses only with unit conversion and measured observables |}\\
\noindent\texttt{| Proposition 37.2 | `normal\_form\_result` | energy-plasma translation can be imported by later papers only through its declared source, receipt, and lane. |}\\
\noindent\texttt{| Protocol 37.3 | `falsifier\_or\_open\_obligation` | laboratory plasma behavior, joule conversion, and universal energy bounds require external data |}\\

\subsection*{Reviewer Claim Ledger}

This ledger expands the formal-claims table into review actions. It is intended to make proof granularity explicit: what is being claimed, what evidence type can support it, what boundary remains, and what the next review action is.

\noindent\texttt{| Formal row | Type | Claim in review terms | Evidence required | Boundary | Next review action |}\\
\noindent\texttt{|---|---|---|---|---|---|}\\
\noindent\texttt{| Theorem 37.1 | external target or comparator | internal traversal cost and carrier thresholds can be translated into plasma/energy hypotheses only with unit conversion and measured observables | dataset, measured value}\\
\noindent\texttt{| Proposition 37.2 | translation grammar | energy-plasma translation can be imported by later papers only through its declared source, receipt, and lane | definitions, normal-form construction, conversion rule, and downs}\\
\noindent\texttt{| Protocol 37.3 | obligation/falsifier | laboratory plasma behavior, joule conversion, and universal energy bounds require external data | explicit missing derivation, adapter, receipt, dataset, comparator, or failed tes}\\

\subsection*{Claim Granularity And Review Table}

The table below separates the claim types so the paper is reviewable without accepting the whole FLCR vocabulary at once.

\noindent\texttt{| Claim type | What this paper may claim | Acceptance test | What is not claimed by that row |}\\
\noindent\texttt{|---|---|---|---|}\\
\noindent\texttt{| Formal-system result | `FLCR-37` defines or uses **Plasma, Energy, And Traversal Calibration** as a typed FLCR object with declared inputs, operations, outputs, and residue. | Definitions, formal claims, construction s}\\
\noindent\texttt{| Computational or receipt-bound result | Enumerations, TarPit runs, CAM/crystal routes, verifier rows, and generated artifacts are claims about the stated finite or executable domain. | A named receipt, validator, manif}\\
\noindent\texttt{| Imported theorem or external definition | The Standard Model or adjacent physics object is closed only as an external target definition when the relevant definition/citation row is present. | The source theorem, defini}\\
\noindent\texttt{| Interpretive bridge or analogy | Analog, workbook, visual, or translation language may explain why the construction is useful. | The analogy preserves the relevant structure and does not promote itself into a theorem. }\\
\noindent\texttt{| Physical or calibration-facing claim | Measured values, physical identity, and predictive accuracy require scale, units, dataset, uncertainty, comparator, and pass/fail criteria. | A dataset, citation, calibration prot}\\
\noindent\texttt{| Open obligation or falsifier | A missing derivation, adapter, receipt, dataset, or failed comparator is a named research channel. | The next binding action and the reason the stronger claim is not closed are stated. | }\\

\subsection*{5. Paper-Specific Development}

1. Identify the local object and its assigned sources for FLCR-37.
2. Separate what is internally receipt-bound from what is citation-bound, CAM-bound, calibration-bound, or still an obligation.
3. State the strongest admissible claim in its current lane.
4. Export only the lane-safe result to downstream papers and preserve the residue.

\subsubsection*{Proof Sketch}

The proof strategy for FLCR-37 is a typed construction argument. The paper names observable as the active object, binds it to the routed legacy and orbital sources, and then allows downstream use only through the declared claim lane. This does not erase stronger ambitions; it keeps them addressable as dependencies, calibration tasks, or falsifiers.

\subsection*{6. Construction}

The construction is intentionally staged. First, the paper names the finite or
typed object that can be inspected directly. Second, it declares the admissible
operations over that object. Third, it records the receipt or source class that
allows another paper to consume the result. Fourth, it names the residue that
must not be silently promoted.

For this paper, the operative object is `Plasma, Energy, And Traversal Calibration`. Its safe consumption
rule is:

\begin{verbatim}
source object -> declared operation -> receipt/lane -> downstream import
\end{verbatim}

The unsafe consumption rule is:

\begin{verbatim}
source phrase -> downstream identity claim
\end{verbatim}

That second pattern is explicitly rejected by FLCR-01 and must be rewritten as
a claim-lane promotion with evidence.

\subsection*{7. Evidence And Receipts}

\noindent\texttt{| Evidence source | What it contributes |}\\
\noindent\texttt{|-----------------|---------------------|}\\
\noindent\texttt{| Legacy paper body | Primary formal and source-integration material assigned by the series map. |}\\
\noindent\texttt{| Internal and pyramid supplements | Cross-paper reasoning windows, deferred back-application slots, and route constraints. |}\\
\noindent\texttt{| NSIT tool closure matrix | Existing verifier/tool families that can close internal or batch claims. |}\\
\noindent\texttt{| CAM/crystal and standards-window evidence | Projected memory, source routing, standards reports, and window/node databases. |}\\

\subsubsection*{Standard Model Definition Dependencies}

This paper consumes the dedicated Standard Model Defining layer. These papers define the external Standard Model or adjacent standard-physics object before this FLCR paper translates it into LCR language.

\noindent\texttt{| SMD paper | Definition surface | Required use |}\\
\noindent\texttt{|---|---|---|}\\
\noindent\texttt{| `SMD-01` | Standard Model Definition Contract And Evidence Boundary | Definition surface that must be cited before this paper promotes the corresponding Standard Model-facing claim. |}\\
\noindent\texttt{| `SMD-05` | QCD Color Dynamics And Hadronic Boundary | Definition surface that must be cited before this paper promotes the corresponding Standard Model-facing claim. |}\\
\noindent\texttt{| `SMD-08` | Renormalization, Effective Field Theory, And Running Parameters | Definition surface that must be cited before this paper promotes the corresponding Standard Model-facing claim. |}\\
\noindent\texttt{| `SMD-11` | GR Continuum Interface And Units-Bearing Calibration | Definition surface that must be cited before this paper promotes the corresponding Standard Model-facing claim. |}\\
\noindent\texttt{| `SMD-12` | Standard Model Comparator And Falsifier Protocol | Definition surface that must be cited before this paper promotes the corresponding Standard Model-facing claim. |}\\

This paper also imports SMB-01, the Standard Model Closure Bridge. SMB-01 supplies the closure-by-lane rule: rows are no longer treated as unresolved by default; they are closed when standard formalism, FLCR proof, CAM/crystal reapplication, normal-form translation, or external calibration satisfies the lane, and they remain obligations only when a required field is missing.

\subsubsection*{Closed Rows Applied In This Paper}

\noindent\texttt{| Row | Closure | Evidence | Boundary |}\\
\noindent\texttt{|---|---|---|---|}\\
\noindent\texttt{| Internal energetic traversal accounting | closed internally | `FLCR-24`, `FLCR-25` | Closes unit-agnostic traversal rows, not measured energy. |}\\
\noindent\texttt{| Plasma/shear carrier horizon | closed at finite carrier scope | `FLCR-25`, `SMD-11` | Closes structural carrier/shear boundary, not full plasma physics. |}\\
\noindent\texttt{| Units-bearing energy/plasma calibration | not closed here | CODATA/NIST, `SMD-08`, `SMD-11` | Requires units, scale, model, and measured/simulated data. |}\\

The paper can now state that its traversal and shear machinery is a real internal accounting system. The open work is not whether anything exists; it is the specific physical calibration from internal row to measured energy/plasma quantity.

\subsubsection*{Active Comparator Rows}

\noindent\texttt{| Comparator | External side | LCR side | Current result | Next required action |}\\
\noindent\texttt{|---|---|---|---|---|}\\
\noindent\texttt{| Energy units | joule/eV/constants | unitless traversal ledger unless calibrated | calibration-bound | Add dimensional analysis and constants. |}\\
\noindent\texttt{| Plasma model | MHD/kinetic/shear model | carrier/shear horizon | structural only | Name model equations and data source. |}\\
\noindent\texttt{| Monster/ceiling boundary | large invariant bound | energy-bound hypothesis | structural ceiling only | Add units-bearing witness before measured-energy claim. |}\\

\subsubsection*{Executable Evidence Bound In This Pass}

This paper now distinguishes three energy/plasma surfaces: internal traversal accounting, carrier/shear verifier routing, and units-bearing physical calibration. That distinction lets the paper use the forge evidence without pretending that a physical plasma experiment has already been run.

\noindent\texttt{| Evidence | Path | Result imported here | Boundary preserved |}\\
\noindent\texttt{|---|---|---|---|}\\
\noindent\texttt{| Energetic traversal verifier route | `local evidence artifact` | Catalogs `python production/formal-papers/CQE-paper-25/verify\_energetic\_traversal.py`; expected status `pass\_with\_open\_obligations`; checks NSLTerm arith}\\
\noindent\texttt{| Z-pinch/shear verifier route | `local evidence artifact` | Catalogs `python production/formal-papers/CQE-paper-26/verify\_zpinch\_shear\_horizon.py`; expected status `pass\_with\_open\_obligations`; checks integer Oloid peri}\\
\noindent\texttt{| Continuum/units definition surface | `SMD-11` | Requires units, equations, constants, uncertainty, and external data for energy, gravity, and plasma claims. | Keeps physical calibration separate from internal traversal}\\

The upgraded claim is: FLCR-37 can state that traversal and shear accounting have named verifier routes and internal closure semantics. It cannot promote those rows into measured plasma, energy, or force claims until units-bearing calibration and external data are supplied.

\subsubsection*{Online Source Rows Applied In This Pass}

\noindent\texttt{| Online source | Claim-lane effect | Boundary retained |}\\
\noindent\texttt{|---|---|---|}\\
\noindent\texttt{| `IRL-PPPL-PLASMA-MAGNETIC` | Supplies real plasma/magnetic-confinement context: charged plasma particles move under magnetic-field constraints. | Does not validate any LCR plasma-energy prediction. |}\\
\noindent\texttt{| `IRL-PPPL-PLASMA-DIAGNOSTICS-2026` | Supplies plasma diagnostic framing: measured plasma claims need particle/field diagnostics. | Internal traversal/shear rows remain unit-agnostic until data is attached. |}\\
\noindent\texttt{| `IRL-CODATA-2022-WALLET` | Supplies constants/unit rows for energy/plasma calibration. | Constants do not close a plasma experiment. |}\\

Online data therefore closes the lack of plasma reference material. The remaining open row is experimental: define device, plasma parameters, diagnostic, units, and pass/fail comparator.

\subsubsection*{Assertive Closure Promotions}

\noindent\texttt{| Row | Closure now allowed | Evidence | Residue moved out of the open row |}\\
\noindent\texttt{|---|---|---|---|}\\
\noindent\texttt{| Traversal accounting route | closed as named internal verifier route | `verify\_energetic\_traversal.py` catalog row | Units-bearing energy experiment. |}\\
\noindent\texttt{| Shear/pinch carrier route | closed as named internal verifier route | `verify\_zpinch\_shear\_horizon.py` catalog row | Physical plasma/Z-pinch validation. |}\\
\noindent\texttt{| Plasma source lane | closed by standard plasma/diagnostic references | PPPL plasma sources, `SMD-11` | Device-specific diagnostic data. |}\\

The non-timid conclusion is: FLCR-37 can close the accounting and source lanes. The only honest open row is measured plasma/energy calibration.

\subsection*{Appendix A. Recursive Capability Reapplication Review}

This review treats the newly admitted capability families as operators. The question is not only what each source says, but what it solves when the source is recursively reapplied through CAM, claim lanes, forge admission, receipt loops, and paper-to-paper routing.

\subsubsection*{Operator Effects}

\noindent\texttt{| Operator | Standalone result | Recursive result | Newly resolved state | Remaining boundary |}\\
\noindent\texttt{|---|---|---|---|---|}\\

\subsubsection*{Combined Recursive Effects}

\noindent\texttt{| Combination | Result | Solves |}\\
\noindent\texttt{|---|---|---|}\\
\noindent\texttt{| COMBO-APPLIED-PHYSICS-BOUNDARY | Applied forge topology supplies the domain boundary, GLM closure supplies bounded mathematical promotions, and address horizons supply finite computational evidence. | This lets later S}\\

The practical consequence is that several prior gaps are now better classified. Some are closed as routing, admission, finite-address, or executable-publication problems. The residues that remain are sharper: exhaustive source intake, physical calibration, full paper-local runner parity, external measurement, and domain-specific validation.

\subsection*{Appendix D. Resolved-State Closure Pass}

This pass removes false restrictions on the paper's claim posture. A row is no longer called open merely because a stronger future claim exists. The satisfied lane is closed now; only the stronger claim that lacks its required adapter, comparator, calibration, or falsifier remains as residue.

\subsubsection*{Closed Now}

\noindent\texttt{| Row | Lane | Resolved state | Remaining boundary |}\\
\noindent\texttt{|---|---|---|---|}\\
\noindent\texttt{| paper lift enumeration | `receipt\_bound\_internal\_result` | 8/8 LCR states succeeded; error walls=0; avg TarPit mass=0.505264. | This closes the paper-local finite lift readout, not every stronger physical interpretatio}\\
\noindent\texttt{| causal ledger | `normal\_form\_result` | The causal/obligation ledger is closed as a dependency-preserving channel, not as an unresolved label. | The family closure is a structural lane; measured physics still requires d}\\
\noindent\texttt{| decade-4 tower | `receipt\_bound\_internal\_result` | The decade tower is resolved with avg TarPit mass=0.509077 and mass sum=40.726174. | The decade total is an internal tower metric, not a standalone universal physical }\\
\noindent\texttt{| family-07 cross-lift comparison | `cam\_crystal\_reapplication\_result` | Family tower binds FLCR-07, FLCR-17, FLCR-27, FLCR-37; avg TarPit mass=0.507568; error walls=0. | Cross-lift agreement strengthens routing and comp}\\
\noindent\texttt{| plasma traversal calibration | `external\_calibration\_result` | Plasma, shear, and traversal source lanes are closed enough to define comparator rows. | Experiment/device claims require diagnostics. |}\\

\subsubsection*{Claim Posture After This Pass}

`FLCR-37` should now be read as a resolved-state contribution for `Plasma, Energy, And Traversal Calibration` at its declared lane. The paper may state the strongest claim supported by these rows directly. It should reserve caution only for a specifically named stronger claim whose required evidence is absent.

\subsection*{8. Dependencies And Downstream Use}

This paper may be imported by later FLCR papers only through the claim lanes
above. The default downstream consumers are:

\noindent{}-- `FLCR-11` through `FLCR-18` for window, receipt, admission, CA, algebraic,\\
  curvature, residue, and exceptional machinery.
\noindent{}-- `FLCR-28` through `FLCR-30` for CAM/crystal, corpus ribbon, and universal\\
  normal-form intake.
\noindent{}-- `FLCR-31` through `FLCR-40` only after the LCR-native result has been cited\\
  and the translation claim has its own evidence lane.

\subsection*{9. Limitations And Falsifiers}

\noindent{}-- laboratory plasma behavior, joule conversion, and universal energy bounds require external data\\
\noindent{}-- External calibration claims require units, datasets, citations, and reproducible data binding.\\
\noindent{}-- A later translation paper may strengthen this result only by adding the missing lane evidence.\\

A falsifier for this paper is any example that satisfies the declared input
conditions while violating the stated construction, receipt, or lane boundary.
An interpretation that merely wants a stronger downstream conclusion is not a
falsifier; it is an obligation routed to a later paper.

\subsection*{10. Publication Readiness Tasks}

\noindent{}-- Validator identities and receipt hashes are recorded in the manifest or receipt appendix.\\
\noindent{}-- External theorems require final citation entries before journal submission.\\
\noindent{}-- Every theorem, proposition, and protocol must retain a claim lane and evidence boundary.\\
\noindent{}-- The Markdown, LaTeX, and PDF forms are generated from the same canonical source.\\

\subsection*{Dual Positioning: Story And Formal Carrier}

This paper must be read in two synchronized ways. Sequentially, it explains why
the next state exists in the corpus story. Formally, it binds that state to
accepted mathematics, IRL formalism, code receipts, validators, CAM/crystal
lookups, and claim-lane boundaries.

Assigned original ribbon role(s): `25`/carrier\_action, `26`/ledger\_action, `29`/window\_action.

\noindent\texttt{| Original slot | Family | Lift | Role | Proof form |}\\
\noindent\texttt{|---|---|---:|---|---|}\\
\noindent\texttt{| `25` | carrier\_action | 2 | order-3 slot-5: carry the state through a path, traversal, or admissible motion | carrier/transducer/path proof |}\\
\noindent\texttt{| `26` | ledger\_action | 2 | order-3 slot-6: bind state into causal, observer, or synchronization ledgers | ledger, graph, and synchronization proof |}\\
\noindent\texttt{| `29` | window\_action | 2 | order-3 slot-9: read the ribbon through windows, meta-readouts, or synthesis bands | window count, source preservation, and meta-ribbon proof |}\\

The formal obligation is to state the strongest valid claim available for this
paper without letting interpretation silently change the addressed object. Any
author interpretation belongs beside the formalism as a declared relabeling,
bridge, or analog, and must be accompanied by the evidence lane that makes the
claim consumable by later papers.

\subsection*{State-Bound Proof Contract}

This paper receives: state emitted by prior slot 24 and reopened at original slot 25 (Energetic Traversal Maps).

It must produce: formal result of original slot 29 plus the same-family lift path toward slot 39.

Semantic transition: carry the state through an admissible path, traversal, transport, or transducer; bind state changes into causal, observer, synchronization, or obligation ledgers; read the ribbon through temporal, observer, Hamiltonian, meta, or synthesis windows.

Accepted formalism targets to bind in the journal rewrite:

\noindent{}-- path composition\\
\noindent{}-- automata and transducers\\
\noindent{}-- transport maps\\
\noindent{}-- invariant preservation along morphisms\\
\noindent{}-- directed graphs and partial orders\\
\noindent{}-- causal/event ledgers\\
\noindent{}-- synchronization protocols\\
\noindent{}-- traceability and provenance invariants\\
\noindent{}-- windowed dynamical systems\\
\noindent{}-- operator/readout notation\\
\noindent{}-- Hamiltonian or energy-style accounting when explicitly bounded\\
\noindent{}-- multi-scale dependency summaries\\

\noindent\texttt{| Slot | Family | Lift | Produced proof form |}\\
\noindent\texttt{|---|---|---:|---|}\\
\noindent\texttt{| `25` | carrier\_action | 2 | carrier/transducer/path proof |}\\
\noindent\texttt{| `26` | ledger\_action | 2 | ledger, graph, and synchronization proof |}\\
\noindent\texttt{| `29` | window\_action | 2 | window count, source preservation, and meta-ribbon proof |}\\

The formal carrier uses these accepted tools while preserving interpretive
claims as explicit relabelings, correspondences, or bridges. Claim promotion is admissible only when a receipt, standard citation,
CAM/crystal reapplication, normal-form result, calibration datum, or falsifier
boundary is attached.

\subsection*{Provenance And Crystal Evidence Integration}

This section integrates prior work, CAM/crystal projections, NSIT tooling,
standards-window evidence, and routed supplements as provenance-bearing
evidence. Source material is not treated as manuscript text; it is bound into
the paper only through claim lanes, receipts, validators, normal forms,
calibration rows, or explicit falsifier boundaries.

\subsubsection*{Expansion Rule For This Slot}

Use past work to bind transport: path, transducer, traversal, carrier medium, invariant, and external calibration boundary.

\subsubsection*{Primary Evidence Inputs}

\noindent\texttt{| Source | Placement reason |}\\
\noindent\texttt{|---|---|}\\
\noindent\texttt{| `25\_Energetic\_Traversal\_Maps.md` | primary assigned source for the paper's formal spine |}\\
\noindent\texttt{| `26\_Z\_Pinch\_and\_Shear\_Horizon.md` | primary assigned source for the paper's formal spine |}\\
\noindent\texttt{| `29\_Monster\_Universal\_Energy\_Bound\_Hypotheses.md` | primary assigned source for the paper's formal spine |}\\
\noindent\texttt{| `supplements/25\_INTERNAL\_REASONING\_SUPPLEMENT.md` | paper-local reasoning supplement contributing to definitions, proof sketch, and limitations |}\\
\noindent\texttt{| `supplements/26\_INTERNAL\_REASONING\_SUPPLEMENT.md` | paper-local reasoning supplement contributing to definitions, proof sketch, and limitations |}\\
\noindent\texttt{| `supplements/29\_INTERNAL\_REASONING\_SUPPLEMENT.md` | paper-local reasoning supplement contributing to definitions, proof sketch, and limitations |}\\
\noindent\texttt{| `virtuous\textbackslash{}25\_Energetic\_Traversal\_Maps\_VIRTUOUS.md` | prior enriched paper body; should be mined for mature claims and proof ordering |}\\
\noindent\texttt{| `virtuous\textbackslash{}26\_Z-Pinch\_and\_Shear\_Horizon\_VIRTUOUS.md` | prior enriched paper body; should be mined for mature claims and proof ordering |}\\
\noindent\texttt{| `virtuous\textbackslash{}29\_Monster\_Universal\_Energy-Bound\_Hypotheses\_VIRTUOUS.md` | prior enriched paper body; should be mined for mature claims and proof ordering |}\\

\subsubsection*{Secondary And Orbital Evidence Inputs}

\noindent\texttt{| Source | Placement reason |}\\
\noindent\texttt{|---|---|}\\
\noindent\texttt{| `supplements/SM\_BRIDGE\_SUPPLEMENT.md` | cross-cutting supplement; paper-relevant rows, receipts, and guard language are bound through evidence lanes |}\\
\noindent\texttt{| `supplements/METAFORGE\_MATERIALS\_SUPPLEMENT.md` | cross-cutting supplement; paper-relevant rows, receipts, and guard language are bound through evidence lanes |}\\
\noindent\texttt{| `supplements/RECEIPT\_VERIFIER\_CATALOG.md` | cross-cutting supplement; paper-relevant rows, receipts, and guard language are bound through evidence lanes |}\\
\noindent\texttt{| `supplements/APPLIED\_FORGES\_WORKBOOK.md` | cross-cutting supplement; paper-relevant rows, receipts, and guard language are bound through evidence lanes |}\\
\noindent\texttt{| `supplements/INTERNAL\_REASONING\_PYRAMID\_INDEX.md` | cross-cutting supplement; paper-relevant rows, receipts, and guard language are bound through evidence lanes |}\\
\noindent\texttt{| `supplements/INTERNAL\_REASONING\_5P\_WINDOW\_SUPPLEMENT.md` | cross-cutting supplement; paper-relevant rows, receipts, and guard language are bound through evidence lanes |}\\
\noindent\texttt{| `supplements/INTERNAL\_REASONING\_7P\_WINDOW\_SUPPLEMENT.md` | cross-cutting supplement; paper-relevant rows, receipts, and guard language are bound through evidence lanes |}\\
\noindent\texttt{| `supplements/INTERNAL\_REASONING\_9P\_WINDOW\_SUPPLEMENT.md` | cross-cutting supplement; paper-relevant rows, receipts, and guard language are bound through evidence lanes |}\\
\noindent\texttt{| `CAM\_CLAIM\_100\_SCORE\_AUDIT.md` | series-wide audit/score/closure evidence; import paper-specific rows and leave global context outside the body |}\\
\noindent\texttt{| `NSIT\_TOOL\_CLOSURE\_MATRIX.md` | series-wide audit/score/closure evidence; import paper-specific rows and leave global context outside the body |}\\
\noindent\texttt{| `NSIT\_REASONED\_CLOSURE\_PASS.md` | series-wide audit/score/closure evidence; import paper-specific rows and leave global context outside the body |}\\
\noindent\texttt{| `ORBITAL\_DATA\_CROSS\_POPULATION\_AUDIT.md` | series-wide audit/score/closure evidence; import paper-specific rows and leave global context outside the body |}\\
\noindent\texttt{| Additional source routes | Complete routing is maintained in the source-placement appendix |}\\

\subsubsection*{Crystal And Standards Evidence To Bind}

\noindent{}-- Paper Reworks crystal projection: 33 paper markdown files, 9 supplements, 5 queues, 6 lattice-forge unification artifacts, 140 faces, 140 vignettes, and 280 CAM nodes.\\
\noindent{}-- Standards-window intake: 192/192 corpus conformance verdicts PASS; 140/142 obligations have candidate routes; 2/4/8/16/32 window lattice is available for cross-paper reads.\\
\noindent{}-- Agent/crystal harvest: agent-generated memory, runtime standards starter, current runtime build, repo-harvest CAM, notebook-derived bundles, and downloaded crystal payloads are orbital evidence surfaces.\\
\noindent{}-- NSIT inventory baseline: 220 validators, 1786 receipt/data artifacts, 1137 formal-paper-like artifacts, 114 supplements, and 86 memory/CAM artifacts.\\

\subsubsection*{Paper-Specific CAM/Score Rows}

\noindent\texttt{| 25 | 88 | internal energy ledger closed | CL verifier: kappa=ln(phi)/16 conservation descent, energy ledger 5/5 | joule conversion and physical units calibration |}\\
\noindent\texttt{| 26 | 80 | carrier-level pinch/shear closed | carrier-level Z-pinch/shear horizon, Paper 14 curvature/repair, Paper 25 energy support | measured plasma/Z-pinch observable and calibration |}\\
\noindent\texttt{| 29 | 90 | internal Monster map closed | CL verifier: 196883=47*59*71, McKay anchors 783/4372, Paper 18 route support | measured universal energy-bound witness and encoder invariance |}\\

\subsubsection*{Paper-Specific NSIT Closure Rows}

No direct NSIT row matched the assigned legacy papers in the first-pass matrix; validator matching by object name remains a receipt-attachment requirement.

\subsubsection*{Source Integration Summary}

The legacy source and internal reasoning supplement are treated as source
evidence rather than manuscript prose. Their contributions enter this paper as
claim-lane entries, receipt or validator references, finite-domain statements,
and limitation boundaries. Speculative or cross-domain interpretations remain
separate from closed local results until an explicit adapter, citation,
normal-form derivation, calibration row, or falsifier is attached.
\subsection*{Claim Binding Queue}

This queue records the evidence discipline for claim strengthening. It separates claims that already have support from residues that remain in falsifier or open-obligation lanes.

\noindent\texttt{| Queue | Lane | Promote now | Bind next | Residue |}\\
\noindent\texttt{|---|---|---|---|---|}\\
\noindent\texttt{| Leech/E8/Golay Operational Lookup | `receipt\_bound\_internal\_result` | Promote no-cost library lookup, code-chain, E8/Niemeier/Leech construction-surface claims when the lattice-forge API or receipt is named. | Attach e}\\
\noindent\texttt{| Symbolic Carrier Versus Physical Carrier | `external\_calibration\_result` | Promote addressable symbolic-carrier and transport claims where the paper names carrier, path, residue, or traversal rules. | Map every physica}\\
\noindent\texttt{| Moonshine/VOA Finite Sector Split | `receipt\_bound\_internal\_result` | Promote finite VOA sector, centroid-chain, and windowed sector claims when the receipt is attached. | Attach centroid/VOA receipt, finite sector cou}\\
\noindent\texttt{| Continuum Edge And Sampling-Stability Bridge | `external\_calibration\_result` | Promote discrete-continuous bridge claims only when the sample rule, norm, error bound, or counterexample policy is stated. | Attach sampli}\\

\subsubsection*{Binding Discipline}

Each queue item is represented as a delta:

\begin{verbatim}
local claim at paper time
-> attached later source/tool/receipt/theorem
-> revised representation
-> invariant boundary and remaining residue
\end{verbatim}

This preserves the sequential story while allowing later tools, crystals, and
standard formalisms to return through the proper lane.

\subsection*{Claim Delta Ledger}

This ledger applies the paper's claim-binding queues as explicit back-application
deltas. It keeps the sequential paper story intact while allowing later
receipts, standard formalisms, CAM/crystal routes, and validators to sharpen the
claim.

\noindent\texttt{| Delta | Local claim at paper time | Later evidence | Revised representation | Boundary kept |}\\
\noindent\texttt{|---|---|---|---|---|}\\
\noindent\texttt{| Leech E8 Golay Lookup | This paper may name lattice closure, E8/Niemeier/Leech surfaces, or terminal lattice addresses only at the scope stated locally. | Lattice-forge code-chain receipts, E8/Niemeier lookup machinery}\\
\noindent\texttt{| Symbolic Vs Physical Carrier | This paper may use carrier language for addressable symbolic transport inside the LCR/CAM system. | Standard physics carrier terminology, local verifier receipts, material/plasma/transpor}\\
\noindent\texttt{| Moonshine Voa Finite Sector | This paper may use moonshine/VOA language for finite sector or windowed representation surfaces only where locally supported. | Centroid/VOA receipt, finite sector chain, lattice/represent}\\
\noindent\texttt{| Continuum Sampling Bridge | This paper may bridge discrete and continuous language only as far as its sampling, projection, or boundary rule supports. | Sampling/interpolation rules, norms, error bounds, boundary condi}\\

The revised representation records the strongest claim supported by the current evidence package. Promotion into theorem or proposition language occurs only when the named evidence is attached in the manifest or workbook replay.

\subsection*{Source-Backed Evidence Summary}

Routed archive and mirror cards are treated as provenance summaries rather than
body text. They identify candidate receipts, validators, theorem registries, or
supplemental source routes. A card strengthens this paper only when its
contribution is bound to a claim lane, evidence object, and boundary.

\noindent\texttt{| Evidence card | Score | Publication contribution | Boundary |}\\
\noindent\texttt{|---|---:|---|---|}\\
\noindent\texttt{| FORMAL: Paper 23 — C-Form: FoldForge Protein Folding Gluon | 89 | **C = the protein fold Gluon** — the contact-map/topo Gluon that transports protein chain fold hypotheses through contact-map and topology receipts. In }\\
\noindent\texttt{| FORMAL: Paper 23 — C-Form: FoldForge Protein Folding Gluon | 89 | **C = the protein fold Gluon** — the contact-map/topo Gluon that transports protein chain fold hypotheses through contact-map and topology receipts. In }\\
\noindent\texttt{| FORMAL: Paper 23 — C-Form: FoldForge Protein Folding Gluon | 89 | **C = the protein fold Gluon** — the contact-map/topo Gluon that transports protein chain fold hypotheses through contact-map and topology receipts. In }\\
\noindent\texttt{| FORMAL: Paper 22 — C-Form: MetaForge Applied Materials Gluon | 85 | **C = the material Gluon** — the applied materials Gluon that transports the morphic Gluon (Paper 21) into physical material candidates. In the lattic}\\
\noindent\texttt{| FORMAL: Paper 22 — C-Form: MetaForge Applied Materials Gluon | 85 | **C = the material Gluon** — the applied materials Gluon that transports the morphic Gluon (Paper 21) into physical material candidates. In the lattic}\\
\noindent\texttt{| Additional evidence cards | 65 total | The complete card inventory is maintained in the archive evidence matrix. | Cards are source-discovery aids until bound to paper-local evidence. |}\\

\subsection*{Journal Apparatus}

\subsubsection*{Article Type And Intended Venue Fit}

This manuscript is written as a formal-methods and mathematical-systems paper
inside the series *Formalizing LCR, Unifying Standard Models*. Its intended
reader is a technical reviewer who needs claim boundaries, reproducibility
routes, and cross-paper dependencies to be visible without relying on private
context.

\subsubsection*{Structured Contribution Statement}

\noindent\texttt{| Item | Statement |}\\
\noindent\texttt{|---|---|}\\
\noindent\texttt{| Publication ID | `FLCR-37` |}\\
\noindent\texttt{| Legacy source slot(s) | `25, 26, 29` |}\\
\noindent\texttt{| Ribbon role | order-3 slot-5: carry the state through a path, traversal, or admissible motion |}\\
\noindent\texttt{| Proof form | carrier/transducer/path proof |}\\
\noindent\texttt{| Received state | state emitted by prior slot 24 and reopened at original slot 25 (Energetic Traversal Maps) |}\\
\noindent\texttt{| Produced state | formal result of original slot 29 plus the same-family lift path toward slot 39 |}\\

\subsubsection*{Claim-Evidence Table}

\noindent\texttt{| Claim | Lane | Evidence to attach | Boundary |}\\
\noindent\texttt{|---|---|---|---|}\\
\noindent\texttt{| Theorem 37.1 | `external\_calibration\_result` | Receipt, source card, validator, citation, or CAM route named in the paper manifest | internal traversal cost and carrier thresholds can be translated into plasma/energy h}\\
\noindent\texttt{| Proposition 37.2 | `normal\_form\_result` | Receipt, source card, validator, citation, or CAM route named in the paper manifest | energy-plasma translation can be imported by later papers only through its declared source}\\
\noindent\texttt{| Protocol 37.3 | `falsifier\_or\_open\_obligation` | Receipt, source card, validator, citation, or CAM route named in the paper manifest | laboratory plasma behavior, joule conversion, and universal energy bounds require e}\\

\subsubsection*{Figure Plan}

\noindent\texttt{| Figure | Purpose | Caption |}\\
\noindent\texttt{|---|---|---|}\\
\noindent\texttt{| FLCR-37-F1 | Carrier path and invariant payload | Schematic showing how `FLCR-37` turns its received state into the produced state while preserving claim lanes and residue boundaries. |}\\
\noindent\texttt{| FLCR-37-F2 | Evidence routing map | Diagram of source papers, archive cards, CAM/crystal routes, validators, and workbook replay paths used by this manuscript. |}\\
\noindent\texttt{| FLCR-37-F3 | Same-family lift context | Diagram placing this paper in the original `00-79` ribbon family and showing predecessor/successor slots. |}\\

\subsubsection*{In-Text Figure FLCR-37-F1: Carrier path and invariant payload}

\begin{verbatim}
received state
  -> Carrier path and invariant payload
  -> formal claim lane
  -> evidence object / receipt / citation
  -> produced state
  -> remaining residue
\end{verbatim}

\subsubsection*{Table Plan}

\noindent\texttt{| Table | Purpose |}\\
\noindent\texttt{|---|---|}\\
\noindent\texttt{| FLCR-37-T1 | Carrier transport table |}\\
\noindent\texttt{| FLCR-37-T2 | Claim-lane/evidence-boundary table |}\\
\noindent\texttt{| FLCR-37-T3 | Archive evidence card and source-backed upgrade table |}\\
\noindent\texttt{| FLCR-37-T4 | Workbook replay and falsifier table |}\\

\subsubsection*{Worked Example}

Carry one state through a path/transducer and show what remains invariant. The example must name the input object, the operation,
the evidence object, the allowed revised claim, and the remaining boundary.

\subsubsection*{Nomenclature And Glossary}

\noindent\texttt{| Term | Paper-local meaning |}\\
\noindent\texttt{|---|---|}\\
\noindent\texttt{| CAM | Defined by this paper as part of the `carrier\_action` proof role; final definition is recorded in the paper-local glossary. |}\\
\noindent\texttt{| carrier | Defined by this paper as part of the `carrier\_action` proof role; final definition is recorded in the paper-local glossary. |}\\
\noindent\texttt{| claim lane | Defined by this paper as part of the `carrier\_action` proof role; final definition is recorded in the paper-local glossary. |}\\
\noindent\texttt{| crystal | Defined by this paper as part of the `carrier\_action` proof role; final definition is recorded in the paper-local glossary. |}\\
\noindent\texttt{| invariant payload | Defined by this paper as part of the `carrier\_action` proof role; final definition is recorded in the paper-local glossary. |}\\
\noindent\texttt{| path composition | Defined by this paper as part of the `carrier\_action` proof role; final definition is recorded in the paper-local glossary. |}\\
\noindent\texttt{| receipt | Defined by this paper as part of the `carrier\_action` proof role; final definition is recorded in the paper-local glossary. |}\\
\noindent\texttt{| transport | Defined by this paper as part of the `carrier\_action` proof role; final definition is recorded in the paper-local glossary. |}\\

\subsubsection*{Appendix FLCR-37-A: Evidence Package}

The appendix records:

1. Legacy source excerpts or source-card references used in the final claim ledger.
2. Receipt and verifier paths, including pass/fail/partial status.
3. CAM/crystal query or projection rows used by the paper, with source identity and boundary.
4. Standards-window and NSIT evidence used for conformance or routing.
5. Falsifier, calibration, or external-data rows that remain unresolved.

\subsubsection*{Data And Code Availability}

The publication package contains the canonical Markdown sources, manifests, source-routing matrices, validation reports, and generated PDF/TeX outputs.

Code and verifier references are cited through manifests, archive evidence cards, receipt catalogs, or workbook replay tables. External datasets or
standards citations must be attached before any calibration-bound claim is
promoted.

\subsubsection*{Reviewer Checklist}

\noindent\texttt{| Check | Reviewer question |}\\
\noindent\texttt{|---|---|}\\
\noindent\texttt{| Claim lane | Does every strong claim declare its lane and evidence type? |}\\
\noindent\texttt{| Boundary | Does the paper preserve what it does not prove? |}\\
\noindent\texttt{| Reproducibility | Is there a receipt, verifier, source card, citation, or replay table for the claim? |}\\
\noindent\texttt{| Cross-paper import | Does the paper cite the prior FLCR/OPR slot before consuming its result? |}\\
\noindent\texttt{| Interpretation | Does the analog layer preserve the addressed state while changing labels/access paths? |}\\

\subsubsection*{Declarations}

No human-subjects, clinical, or animal-research protocol is asserted by this
paper. External empirical claims require separate dataset, calibration, or
domain-validation attachments. Competing-interest and funding statements are recorded in the submission metadata.

\subsection*{11. Publication Integration Addendum}

\subsubsection*{11.1 Role In The Series}

`FLCR-37` (`Plasma, Energy, And Traversal Calibration`) occupies the `carrier\_action` role at lift depth `2`.
The paper receives state emitted by prior slot 24 and reopened at original slot 25 (Energetic Traversal Maps). It produces formal result of original slot 29 plus the same-family lift path toward slot 39. The state transition is:
carry the state through an admissible path, traversal, transport, or transducer; bind state changes into causal, observer, synchronization, or obligation ledgers; read the ribbon through temporal, observer, Hamiltonian, meta, or synthesis windows.

The paper-local proof form is:

\begin{verbatim}
carrier/transducer/path proof
\end{verbatim}

The integration result is a state-preserving carrier or transducer route with recorded loss/residue. This addendum records how that result is
consumed by the publication series without relying on editorial instructions or
private working context.

\subsubsection*{11.2 Formal Carrier}

\noindent\texttt{| Formal carrier | Function |}\\
\noindent\texttt{|---|---|}\\
\noindent\texttt{| path composition | Formal carrier for the paper-local result. |}\\
\noindent\texttt{| automata and transducers | Formal carrier for the paper-local result. |}\\
\noindent\texttt{| transport maps | Formal carrier for the paper-local result. |}\\
\noindent\texttt{| invariant preservation along morphisms | Formal carrier for the paper-local result. |}\\
\noindent\texttt{| directed graphs and partial orders | Formal carrier for the paper-local result. |}\\
\noindent\texttt{| causal/event ledgers | Formal carrier for the paper-local result. |}\\
\noindent\texttt{| synchronization protocols | Formal carrier for the paper-local result. |}\\
\noindent\texttt{| traceability and provenance invariants | Formal carrier for the paper-local result. |}\\
\noindent\texttt{| windowed dynamical systems | Formal carrier for the paper-local result. |}\\
\noindent\texttt{| operator/readout notation | Formal carrier for the paper-local result. |}\\
\noindent\texttt{| Hamiltonian or energy-style accounting when explicitly bounded | Formal carrier for the paper-local result. |}\\
\noindent\texttt{| multi-scale dependency summaries | Formal carrier for the paper-local result. |}\\

The LCR, CAM, crystal, forge, and analog vocabulary functions as a typed access
layer over these carriers. A relabeling is admissible only when the addressed
object, evidence lane, boundary, and downstream import rule remain attached.

\subsubsection*{11.3 Claim Ledger}

\noindent\texttt{| Claim | Lane | Statement |}\\
\noindent\texttt{|---|---|---|}\\
\noindent\texttt{| Theorem 37.1 | `external\_calibration\_result` | internal traversal cost and carrier thresholds can be translated into plasma/energy hypotheses only with unit conversion and measured observables |}\\
\noindent\texttt{| Proposition 37.2 | `normal\_form\_result` | energy-plasma translation can be imported by later papers only through its declared source, receipt, and lane. |}\\
\noindent\texttt{| Protocol 37.3 | `falsifier\_or\_open\_obligation` | laboratory plasma behavior, joule conversion, and universal energy bounds require external data |}\\

These claims are consumed at their stated lane. A stronger downstream statement
creates a new claim envelope with its own evidence object.

\subsubsection*{11.4 Evidence Package}

\noindent\texttt{| Evidence class | Routed items | Status |}\\
\noindent\texttt{|---|---|---|}\\
\noindent\texttt{| Legacy sources | `25\_Energetic\_Traversal\_Maps.md`, `26\_Z\_Pinch\_and\_Shear\_Horizon.md`, `29\_Monster\_Universal\_Energy\_Bound\_Hypotheses.md` | routed evidence |}\\
\noindent\texttt{| Evidence inputs | `CAM\_CLAIM\_100\_SCORE\_AUDIT.md`, `NSIT\_TOOL\_CLOSURE\_MATRIX.md`, `NSIT\_REASONED\_CLOSURE\_PASS.md`, `PAPER\_REWORKS\_CRYSTAL\_PROJECTION.json` | routed evidence |}\\
\noindent\texttt{| CAM/crystal sources | `PAPER\_REWORKS\_CRYSTAL\_PROJECTION.json`, `AGENT\_CRYSTAL\_WORK\_HARVEST.json`, `runtime standards artifact`, `notebook-derived source bundle` | routed evidence |}\\
\noindent\texttt{| Standards-window inputs | `window\_summary.json`, `node DBs`, `window DBs` | routed evidence |}\\
\noindent\texttt{| Receipts | external requirement | not attached in this source package |}\\
\noindent\texttt{| Validators | external requirement | not attached in this source package |}\\

Standards-window, runtime, and notebook-derived artifacts are treated
as evidence-routing surfaces. They support source discovery, conformance
checking, and reapplication, but they do not replace citations, receipts,
normal-form derivations, calibration rows, or falsifiers.

\subsubsection*{11.5 Results Representation}

The paper's results section is organized around five publication objects:

1. the exact object acted on at this slot;
2. the admissible operation or transformation;
3. the receipt, enumeration, derivation, lookup, or external theorem supporting
   the operation;
4. the residue or boundary left after the result;
5. the downstream import rule.

\subsubsection*{11.6 Figures And Tables}

\noindent\texttt{| Figure | Role |}\\
\noindent\texttt{|---|---|}\\
\noindent\texttt{| FLCR-37-F1: Carrier path and invariant payload | Visualizes the proof object, routing, or boundary. |}\\
\noindent\texttt{| FLCR-37-F2: Evidence routing map | Visualizes the proof object, routing, or boundary. |}\\
\noindent\texttt{| FLCR-37-F3: Same-family lift context | Visualizes the proof object, routing, or boundary. |}\\

\noindent\texttt{| Table | Role |}\\
\noindent\texttt{|---|---|}\\
\noindent\texttt{| FLCR-37-T1: Carrier transport table | Records claim lanes, evidence, sources, or falsifiers. |}\\
\noindent\texttt{| FLCR-37-T2: Claim-lane/evidence-boundary table | Records claim lanes, evidence, sources, or falsifiers. |}\\
\noindent\texttt{| FLCR-37-T3: Archive evidence card and source-backed upgrade table | Records claim lanes, evidence, sources, or falsifiers. |}\\
\noindent\texttt{| FLCR-37-T4: Workbook replay and falsifier table | Records claim lanes, evidence, sources, or falsifiers. |}\\

\subsubsection*{11.7 Limits And Falsifiers}

The paper proves only the object and domain declared in its claim ledger.
Physical, Standard Model, observer, calibration, or synthesis claims require
the additional lane evidence stated by their destination papers. CAM/crystal
and forge language is operational only when represented as an address, lookup,
receipt, validator, or reapplication path.
\end{document}
